%% -------------------------------------------------------------------
\subsection{\vrev{이론적 기여}}
\label{subsec:ch5_theoretical}
%% -------------------------------------------------------------------

본 연구의 이론적 기여는 네 가지 차원에서 정리된다.

\chg{첫째, 개별 시그모이드 바운딩에 의한 물리적 경계 보장이다.}
앞서 비교한 선형 RTO 모델의 물리적 경계 위반 문제를 개별 시그모이드
바운딩 함수로 해결하였다. $E_{O,\text{bnd}} \in (-R_0, 0]$과
$E_{U,\text{bnd}} \in [0, R_{\max} - R_0)$의 범위가 해석적으로 보장되어,
$\DRTO \in (0, R_{\max})$가 입력값에 무관하게 성립한다.
이는 기존 BCM 문헌에서 제시되지 않았던 수학적 보장으로서,
\chg{복구목표시간}의 음수 산출이나 MTPD 초과라는 구조적 문제를
근원적으로 해결한다.

\chg{둘째, 이중채널 감쇠 효과의 발견과 정량화이다.}
P85.8 교차점을 기준으로 $k_O$ 계수가 $0.52$배로 감쇠되는
이중채널 현상의 발견은 BCM 이론에 새로운 통찰을 제공한다.
이 발견은 극단 불확실성과 관측성이 비선형적으로 상호 작용하는
메커니즘을 규명하며, 시그모이드 비선형성에 의한 이론적 설명을
제시한다.

\chg{셋째, \citet{lee2026drto}과 본 연구의 일관적 이론 체계 구축이다.}
선행 연구\citep{lee2026drto}의 선형 근사$\to$시그모이드 바운딩$\to$$\kappa$ 최적화에서
본 연구의 이중채널 확장(C2, C3)까지의 단계적 확장이
\chg{수식·데이터·결론}의 교차 정합성을 유지하며 수행되었다.
선형 모델이 시그모이드 바운딩의 1차 근사임을 보이는 관계
($|z| \ll 1$에서 $S(z) \approx z$)는 이론적 확장의
내적 일관성을 수학적으로 보장한다.

  넷째, 불확실성의 확률적 정량화이다. 복구 과정의 불확실성을 일상적 변동을 나타내는 가우시안과 극단 사건을 나타내는 파레토의 두 분포로 정량화함으로써, 관측성 단일 차원에 머물던 기존 모델의 한계를 보완하고 극단 영역의 꼬리 위험을 복구목표시간 산정에 반영하였다.

